\documentclass[11pt,a4paper]{article}
\usepackage[utf8]{inputenc}
\usepackage[margin=1in]{geometry}
\usepackage{amsmath,amssymb,amsfonts}
\usepackage{booktabs}
\usepackage{hyperref}
\usepackage{xcolor}
\usepackage{listings}
\usepackage{microtype}

\hypersetup{
    colorlinks=true,
    linkcolor=blue!70!black,
    citecolor=blue!70!black,
    urlcolor=blue!70!black
}

\title{\textbf{Universal \LaTeX{} Format Standardization Across Project World Model (PWM) Documentation}\\[0.5em]
\large Verification Report \& Formal Mathematical Specification}
\author{\textbf{Petri Paananen}\\
JAMK University of Applied Sciences\\
\texttt{NBN:fi:amk-2026052919003}}
\date{October 17, 2026}

\begin{document}
\maketitle

\begin{abstract}
This document details the universal standardization of mathematical notation, vector space definitions, quasimetric loss functions, and financial metrics across the entire Project World Model (PWM) repository documentation. All pseudo-mathematical expressions, Unicode approximations, and plain-text fraction slashes have been migrated to rigorous \LaTeX{} mathematics. Furthermore, a publication-grade master specification paper (\texttt{Docs/Project\_World\_Model\_Specification.tex}) has been established as the canonical theoretical blueprint for the PWM Causal Digital Twin (CDT) framework.
\end{abstract}

\tableofcontents
\newpage

\section{Introduction \& Architectural Standards}

As formulated in the Master's Thesis research (\textit{Itseohjautuvat työnkulut videopeliteollisuudessa}, Paananen, 2026), the Project World Model (PWM) operates as an L3 Causal Digital Twin (CDT). The system resolves the \textit{Agility Paradox} by shifting production management from reactive issue tracking to proactive, counterfactual simulation in latent space.

To ensure peer-reviewed publication rigor, Section~5 of \texttt{Docs/PLANNING\_DOCUMENT\_STANDARDS.md} mandates that all repository documentation, planning blueprints, and technical whitepapers use formal \LaTeX{} mathematical notation. The canonical notations enforced across the project are summarized in Table~\ref{tab:notation_standards}.

\begin{table}[h!]
\centering
\small
\begin{tabular}{lll}
\toprule
\textbf{Mathematical Concept} & \textbf{Standard \LaTeX{} Notation} & \textbf{Vector / Space Definition} \\
\midrule
Full Latent State Vector & $\mathbf{s} \in \mathbb{R}^{64}$ & Unified latent representation space \\
Factor Coordinates & $\mathbf{c} \in \mathbb{R}^{4 \times 16}$ & $K=4$ orthogonal predictive subspaces \\
Orthonormal Basis Operator & $W \in \mathbb{R}^{64 \times 64}, \quad \|W W^T - I\|_F < 10^{-6}$ & Gram matrix isometry condition \\
Factor Latent Energy & $E_k(\mathbf{c}_k, \mathbf{c}_k') = \|\mathbf{c}_k' - \mathbf{c}_k\|_2$ & Subspace perturbation metric \\
Quasimetric Latent Energy & $E(\mathbf{s}, \mathbf{s}') = \max(0, \Phi(\mathbf{s}') - \Phi(\mathbf{s})) + \epsilon \|\mathbf{s}' - \mathbf{s}\|_2$ & Directed causal distance \\
Risk Tail Bound & $\text{CVaR}_{0.95}(\pi) = \mathbb{E}[R \mid R \ge \text{VaR}_{0.95}(R)]$ & Conditional Value-at-Risk at 95\% \\
Compute-to-Rework Ratio & $\text{CRR} = \frac{C_{\text{compute}}}{V_{\text{rework}}} < 0.10$ & Cognitive FinOps efficiency bound \\
Queueing Invariant & $L = \lambda W$ & Little's Law for task pipelines \\
System Identification & $\mathbf{z}_{\text{env}} = (\gamma_{\text{rev}}, \kappa_{\text{coupling}}, \tau_{\text{latency}}, \sigma_{\text{exo}})^T$ & 4-dimensional domain embedding \\
\bottomrule
\end{tabular}
\caption{Canonical \LaTeX{} notation standards established in \texttt{Docs/PLANNING\_DOCUMENT\_STANDARDS.md}.}
\label{tab:notation_standards}
\end{table}

\section{Document-by-Document Modifications}

\subsection{Core Architecture Documents}

\paragraph{README.md}
Standardized all mathematical formulations across Section~1 (Orthogonal Predictive Factorization), Section~2 (Recursive Self-Improvement), Section~5 (Flow Metrics), Section~6 (Domain-Agnostic CDT), and the CRR Economic Model. Upgraded vector coordinates $\mathbf{s} \in \mathbb{R}^{64}$, Gram condition $\|W W^T - I\|_F < 10^{-6}$, and Quasimetric Latent Energy:
\begin{equation}
E(\mathbf{s}, \mathbf{s}') = \max(0, \Phi(\mathbf{s}') - \Phi(\mathbf{s})) + \epsilon \|\mathbf{s}' - \mathbf{s}\|_2, \quad \epsilon = 0.05
\end{equation}

\paragraph{DTO\_LEGEND.md}
Added Section~12 (\textit{Spatial Physics \& Domain-Agnostic Telemetry Mathematics}). Formatted surface alignment elevation ($y \in [0.005, 0.01]\,\text{m}$), exclusion clearance radii ($r_{\text{fountain}} = 2.4\,\text{m}$, $r_{\text{tree}} = 3.0\,\text{m}$, $r_{\text{agent}} = 1.0\,\text{m}$), and dynamic low-reversibility boulder scaling:
\begin{equation}
\text{scale} = s_0 \cdot \left[1.0 + (1.0 - \gamma_{\text{rev}}) \cdot 0.45\right], \quad w_{\text{col}} = 0.55 \cdot \left(\frac{\text{scale}}{s_0}\right)\,\text{m}
\end{equation}

\subsection{Research Whitepapers}

\paragraph{Docs/JEPA\_Anything\_Review\_and\_PWM\_Integration\_Plan.md}
Converted all Unicode and ASCII representations into formal \LaTeX{}:
\begin{itemize}
    \item State representation mapping: $\mathbf{s} \in \mathbb{R}^d \iff \mathbf{c} \in \mathbb{R}^{K \times r}$.
    \item Dedicated factor predictors: $\hat{\mathbf{c}}_{k, t+1} = f_{\theta_k}(\mathbf{c}_{k, t}, a_t)$.
    \item Isomeric synthesis: $\mathbf{s} = \text{compose}(\mathbf{c}) = W^T \operatorname{vec}(\mathbf{c})$ subject to $W W^T = I$.
    \item Factor risk attribution: $\text{Attr}_k = \frac{E_k}{\sum_{j=1}^K E_j}$.
    \item Multi-factor risk function:
    \begin{equation}
    \text{Risk} = \min\left(\max\left(0.45 E_{\text{topo}} + 0.30 E_{\text{arch}} + 0.15 E_{\text{velo}} + 0.10 E_{\text{coord}} + \text{modifier}, 0.05\right), 0.98\right)
    \end{equation}
\end{itemize}

\paragraph{Docs/RSI\_Anything\_Review\_and\_PWM\_Benchmark.md}
Upgraded the 6-stage closed-loop recursive self-improvement cycle, candidate action sets $\mathcal{A} = \{a_1, a_2, \dots, a_n\}$, and policy pre-screening optimization:
\begin{equation}
\pi^* = \arg\min_{\pi \in \Pi} \left[0.45 \cdot \text{Risk}(\pi) + 0.35 \cdot \text{CVaR}_{0.95}(\pi) + 0.20 \cdot E_{\text{quasi}}(\pi)\right]
\end{equation}

\subsection{Academic Thesis \& Economic Feasibility Studies}

\paragraph{Docs/Petri\_Paananen\_thesis.md}
Upgraded the core Compute-to-Rework Ratio (CRR) formulation in Section~5.8.2 from inline ASCII text into centered display \LaTeX{}:
\begin{equation}
\text{CRR} = \frac{\text{Simulaation inferenssikustannus (€)}}{\text{Vältetyn korjaustyön arvo (€)}}
\end{equation}
The empirical scenario demonstrated in the thesis ($\text{CRR} \approx 0{,}008$) achieves a $125\times$ cost savings multiplier over manual late-stage developer intervention.

\paragraph{Docs/Market Analysis for PWM Framework.md}
Formatted the quantitative enterprise deployment model:
\begin{equation}
\text{CRR} = \frac{\$10}{\$1{,}200} \approx 0{,}0083 \quad (\approx 120\times\ \text{ROI})
\end{equation}
with target operating threshold $\text{CRR} < 0.10$.

\paragraph{Docs/Grand\_Plan.md \& Docs/Grand\_Plan\_Review.md}
Replaced code block pseudo-formulas with display \LaTeX{} and standardized threshold bounds: $\text{CRR} < 0.01$ (Exceptional), $\text{CRR} < 0.10$ (Target), $\text{CRR} < 0.50$ (Acceptable), $\text{CRR} \ge 1.00$ (Danger / Compute Runaway).

\section{Master \LaTeX{} Specification Document}

A standalone, publication-grade academic master paper has been authored at:
\begin{center}
\texttt{Docs/Project\_World\_Model\_Specification.tex}
\end{center}
The document contains:
\begin{enumerate}
    \item Comprehensive 5-layer system architecture with formal algebraic operators.
    \item Orthogonal Predictive Factorization (OPF) theorems and Gram isometry proofs.
    \item Closed-loop Recursive Self-Improvement (RSI) algorithms and pre-screening objectives.
    \item Complete empirical benchmarks across software engineering, game production, and construction domains.
    \item Master bibliography citing foundational works including LeCun (2022), Pearl (2019), and Petri Paananen's Master's Thesis (JAMK, 2026).
\end{enumerate}

\section{Verification \& Quality Assurance}

All local test suites and frontend build targets were executed to verify zero regression:
\begin{itemize}
    \item \textbf{Backend Pytest Suite}: 97 passed in 18.38s (\texttt{.venv\textbackslash Scripts\textbackslash python.exe -m pytest}).
    \item \textbf{Visualizer Frontend}: 0 errors in 1.42s build (\texttt{npm run build} via Vite/Rolldown).
    \item \textbf{Mathematical Syntax Verification}: 100\% matched brackets, valid environments, and zero unescaped pseudo-math delimiters.
\end{itemize}

\end{document}
